% Part: methods
% Chapter: proofs
% Section: example-2

\documentclass[../../../include/open-logic-section]{subfiles}

\begin{document}

\olfileid[hi]{mth}{prf}{ex2}

\olsection{एक और उदाहरण}

\begin{prop}
यदि $A \subseteq C$, तो $A \cup (C \setminus A) = C$।
\end{prop}  

\begin{proof}
  मान लें कि $A \subseteq C$। हम दिखाना चाहते हैं कि $A \cup (C
  \setminus A) = C$।
  \begin{quote}
    हम यह देखकर आरंभ करते हैं कि यह एक सशर्त कथन है। इसका सार्वत्रिक
    परिमाणीकरण अंतर्निहित है: प्रतिज्ञप्ति सभी समुच्चयों $A$ और $C$ के लिए
    सत्य है। अतः $A$ और $C$ मनमाने समुच्चयों के चर हैं। ऐसा कथन सिद्ध करने
    के लिए हम पूर्ववर्ती को मानकर अनुवर्ती सिद्ध करते हैं।

    हम मान्यता $A \subseteq C$ का उपयोग करके आगे बढ़ते हैं। परिभाषा
    $\subseteq$ खोलें: मान्यता का अर्थ है कि $A$ के सभी !!{element}, $C$
    के भी !!{element} हैं। इसे लिख लें---यह एक महत्त्वपूर्ण तथ्य है जिसका
    उपयोग हम पूरे प्रमाण में करेंगे।
  \end{quote}
  परिभाषा~$\subseteq$ से, चूँकि $A \subseteq C$, इसलिए सभी $z$ के लिए,
  यदि $z \in A$, तो $z \in C$।
  \begin{quote}
    हमने मान्यता में दी गई सभी परिभाषाएँ खोल ली हैं। अब हम निष्कर्ष पर बढ़
    सकते हैं। हम दिखाना चाहते हैं कि $A \cup (C \setminus A) = C$, अतः
    पिछले उदाहरण की तरह प्रमाण व्यवस्थित करते हैं: दिखाते हैं कि $A \cup
    (C \setminus A)$ का प्रत्येक !!{element}, $C$ का भी !!a{element} है,
    और इसके विपरीत $C$ का प्रत्येक !!{element}, $A \cup (C \setminus
    A)$ का !!a{element} है। इसे संक्षेप में लिख सकते हैं: $A \cup (C
    \setminus A) \subseteq C$ और $C \subseteq A \cup (C \setminus A)$।
    (यहाँ हम परिभाषा खोलने के विपरीत कर रहे हैं, किंतु इससे प्रमाण पढ़ना कुछ
    सरल हो जाता है।) चूँकि यह एक संयोजन है, हमें दोनों भाग सिद्ध करने हैं।
    पहला भाग दिखाने के लिए, अर्थात् यह कि $A \cup (C \setminus A)$ का
    प्रत्येक !!{element}, $C$ का भी !!a{element} है, हम किसी मनमाने~$z$
    के लिए $z \in A \cup (C \setminus A)$ मानते हैं और दिखाते हैं कि $z
    \in C$। $\cup$ की परिभाषा से, $z \in A \cup (C \setminus A)$ से हम
    निष्कर्ष निकाल सकते हैं कि $z \in A$ या $z \in C \setminus A$। अब
    आपको इसकी रीति समझ आने लगी होगी।
    \end{quote}
  $A \cup (C \setminus A) = C$ तभी और केवल तभी जब $A \cup (C
  \setminus A) \subseteq C$ और $C \subseteq (A \cup (C \setminus A)$।
  पहले हम सिद्ध करते हैं कि $A \cup (C \setminus A) \subseteq C$। मान
  लें $z \in A \cup (C \setminus A)$। अतः या तो $z \in A$ या $z \in
  (C \setminus A)$।
  \begin{quote}
    हम एक वियोजन पर पहुँचे हैं और उससे $z \in C$ सिद्ध करना चाहते हैं। यह
    हम स्थितियों द्वारा प्रमाण से करते हैं।
  \end{quote}
  स्थिति 1: $z \in A$। चूँकि सभी $z$ के लिए, यदि $z \in A$, तो $z \in
  C$, इसलिए हमारे पास $z \in C$ है।
  \begin{quote}
    यहाँ हमने पहले लिखा हुआ वह तथ्य प्रयोग किया है जो प्रतिज्ञप्ति की
    परिकल्पना $A \subseteq C$ से निष्पन्न हुआ था। पहली स्थिति पूरी है और
    अब हम दूसरी स्थिति $z \in (C \setminus A)$ लेते हैं। याद करें कि $C
    \setminus A$ दोनों समुच्चयों का \emph{अंतर} दर्शाता है, अर्थात् $C$ के
    उन सभी !!{element} का समुच्चय जो $A$ के !!{element} नहीं हैं। किंतु
    $C$ का कोई भी !!{element} जो $A$ में नहीं है, विशेषतः $C$ का
    !!a{element} है।
  \end{quote}
  स्थिति 2: $z \in (C \setminus A)$। इसका अर्थ है कि $z \in C$ और $z
  \notin A$। अतः विशेषतः $z \in C$।
  \begin{quote}
    बढ़िया, हमने पहली दिशा सिद्ध कर ली है। अब दूसरी दिशा लें। यहाँ हम सिद्ध
    करते हैं कि $C \subseteq A \cup (C \setminus A)$। अतः हम $z \in C$
    मानकर सिद्ध करते हैं कि $z \in A \cup (C \setminus A)$।
  \end{quote}
  अब मान लें $z \in C$। हम दिखाना चाहते हैं कि $z \in A$ या $z \in C
  \setminus A$।
  \begin{quote}
    चूँकि $A$ के सभी !!{element}, $C$ के भी !!{element} हैं और $C
    \setminus A$ उन सभी वस्तुओं का समुच्चय है जो $C$ के !!{element} हैं
    पर $A$ के नहीं, इसलिए निष्पन्न होता है कि $z$ या तो $A$ में है या $C
    \setminus A$ में। यदि आप पहले से नहीं जानते कि परिणाम सत्य क्यों है, तो
    यह कुछ अस्पष्ट हो सकता है। इसे चरण-दर-चरण सिद्ध करना बेहतर होगा। एक सरल
    तथ्य उपयोगी होगा जिसे हम बिना प्रमाण के कह सकते हैं: $z \in A$ या $z
    \notin A$। इसे ``अपवर्जित मध्य का सिद्धांत'' कहते हैं: किसी भी कथन~$p$
    के लिए, या तो $p$ सत्य है या उसका निषेध सत्य है। (यहाँ $p$ कथन $z \in
    A$ है।) चूँकि यह एक वियोजन है, हम फिर स्थितियों द्वारा प्रमाण कर सकते
    हैं।
  \end{quote}
  या तो $z \in A$ या $z \notin A$। पहली स्थिति में $z \in A \cup (C
  \setminus A)$। दूसरी स्थिति में $z \in C$ और $z \notin A$, अतः $z
  \in C \setminus A$। किंतु तब $z \in A \cup (C \setminus A)$।
  \begin{quote}
    हमारा प्रमाण पूरा हुआ: हमने दिखाया है कि $A \cup (C \setminus A) = C$।
  \end{quote}
\end{proof}

\end{document}
